% 3장/4장 분량 확인용 래퍼.
% 컴파일: (paper-icaif26/ 에서)
%   xelatex -output-directory=build build/build-ko.tex   ... 은 경로가 꼬이므로
%   아래처럼 paper-icaif26/ 에서 실행한다:
%   cp build/build-ko.tex . && xelatex build-ko && bibtex build-ko && xelatex build-ko && xelatex build-ko

\documentclass[sigconf,nonacm]{acmart}

\IfFileExists{kotex.sty}{\usepackage{kotex}}{%
  \usepackage{fontspec}%
  \IfFontExistsTF{Noto Serif CJK KR}{\setmainfont{Noto Serif CJK KR}}{%
    \setmainfont{Noto Serif CJK JP}}%
  \XeTeXlinebreaklocale "ko"%
  \XeTeXlinebreakskip = 0pt plus 0.1em%
}

\usepackage{booktabs}
\usepackage{graphicx}
\usepackage{amsmath}
\usepackage{array}

\settopmatter{printacmref=false,printfolios=true}
\setcopyright{none}

\makeatletter
\let\proposition\relax \let\endproposition\relax
\makeatother
\theoremstyle{acmplain}
\newtheorem{proposition}{명제}[section]

\AtBeginDocument{%
  \renewcommand{\figurename}{그림}%
  \renewcommand{\tablename}{표}%
}

\begin{document}

\title{[분량 확인] SPR-IRL 방법과 실험}
\author{저자명}
\affiliation{\institution{소속기관}\city{서울}\country{대한민국}}
\email{author@example.com}
\renewcommand{\shortauthors}{저자명}
\maketitle

% =====================================================================
%  3장 방법 / 4장 실험 (한국어) — ICAIF 8쪽 제한 압축본
%
%  삽입:  % =====================================================================
%  3장 방법 / 4장 실험 (한국어) — ICAIF 8쪽 제한 압축본
%
%  삽입:  % =====================================================================
%  3장 방법 / 4장 실험 (한국어) — ICAIF 8쪽 제한 압축본
%
%  삽입:  \input{sections/03-04-ko}
%  패키지: kotex, booktabs, graphicx, amsmath, array
%  그림:  figures/  (make_paper_figures.py 산출물 중 3개만 사용)
%  참고문헌: paper-ko.bib + paper-refs-additional.bib
%
%  분량 설계: 3장 약 1.6쪽 + 4장 약 2.4쪽 = 약 4쪽.
%  1~2장에 3쪽, 5장+참고문헌에 1쪽을 남긴다.
% =====================================================================

\section{방법: SPR-IRL}
\label{sec:method}

본 절은 대칭 선호복원(Symmetric Preference Recovery, 이하 SPR-IRL) 절차를
정의한다. SPR-IRL은 모형 하나가 아니라 (i) 공유 특징 위의 대칭 추정, (ii) 표본외
안정성 평가, (iii) 사전 부호 기대와의 대조로 이루어진 절차다.

\subsection{연속 행동과 대칭 추정}
\label{sec:setup}

투자자 유형을 $\mathcal{I}=\{\text{외국인},\text{기관},\text{개인}\}$으로 둔다.
유형 $i$의 $t$일 총매수ㆍ총매도금액을 $V^{buy}_{i,t}$, $V^{sell}_{i,t}$라 하면
행동을
\begin{equation}
u_{i,t}=\frac{V^{buy}_{i,t}-V^{sell}_{i,t}}
             {V^{buy}_{i,t}+V^{sell}_{i,t}}\in[-1,1]
\label{eq:action}
\end{equation}
로 정의한다. 분모가 시장 전체가 아니라 해당 유형의 총거래금액이므로 $u_{i,t}$는
유형별 거래 규모에 영향받지 않는 공통 척도가 되며, 이진 라벨과 달리 강도 정보를
보존한다.

세 유형은 동일한 상태 특징 $x_{i,t}\in\mathbb{R}^{3}$, 동일한 컨텍스트
$C_t\in\mathbb{R}^{2}$, 동일한 모형 형태를 부여받고 파라미터
$\theta_i=(\beta_i,B_i,\alpha_i)$만 각자의 데이터로 분리 학습한다. 유형별로 다른
특징을 미리 배정하지 않으므로, 결과의 유형 간 차이는 설계가 아니라 데이터에서
학습된 것이며 이것이 계수의 직접 비교를 정당화한다.

학습 표본은 $(x_{i,t},C_t,a_{i,t+1})$, $a_{i,t+1}=u_{i,t+1}$이다. $t$일 상태가
$t+1$일 행동을 설명하므로 미래 거래는 특징에 들어가지 않는다.

\subsection{보상 특징과 시장 컨텍스트}
\label{sec:features}

세 특징과 두 컨텍스트는 모두 각 분할의 학습 인덱스에서만 표준화하므로, 이후 계수는
1표준편차 변화의 효과로 읽는다.

\textbf{(F1) 중기 모멘텀} $x^{mom}_{t}=\log(P_t/P_{t-20})$. 20거래일은 해석을 위한
사전 지정값이며 자료에서 최적화하지 않았다. 양수는 추세추종, 음수는 역추세를 뜻하고
부호 기대는 유형별 모멘텀 거래 문헌~\cite{grinblatt2000investment,choe1999foreign}에서
온다.

\textbf{(F2) 군집행동} 유형 $i$를 제외한 두 유형의 전일 행동 횡단면 평균
\begin{equation}
x^{herd}_{i,t}=\tfrac{1}{2}\left(u_{j,t-1}+u_{k,t-1}\right),\qquad j,k\neq i.
\label{eq:herd}
\end{equation}
1거래일 지연은 정보 제약이 아니라, 동일 거래일 유형 간 수급의 회계적 상쇄를 예측
변수에서 배제하기 위한 설계 선택이다(\S\ref{sec:beta}). 부호 기대는 기관 군집행동
문헌~\cite{lakonishok1992impact}에서 오나, 해당 문헌은 기관 \emph{내부}의 군집을,
식~\eqref{eq:herd}은 \emph{유형 간} 관계를 측정한다.

\textbf{(F3) 평균단가 대비 손실구간} 평균단가 대용치 $\bar P_{i,t}$와 당일 총거래
VWAP $P^{vwap}_{i,t}$로
\begin{equation}
x^{uw}_{i,t}=\max\!\left(
\frac{\bar P_{i,t}-P^{vwap}_{i,t}}{\bar P_{i,t}},\,0\right).
\label{eq:underwater}
\end{equation}
0에서 절단하므로 이 변수는 준거점 의존적 설계다. 평균단가는 총매수ㆍ총매도 수량과
금액으로 유효 보유량을 재귀 갱신해 구하되 전일 보유량에 망각계수 $\rho=0.98$을
곱한다(반감기 약 34거래일). 부호 기대는 처분효과
문헌~\cite{shefrin1985disposition,odean1998investors}에서 온다. 다만 이 대용치는
집계 거래로 만든 근사치이므로 $x^{uw}$의 계수를 처분효과의 직접 검정으로 읽을 수
없다.

\textbf{컨텍스트}는 KOSPI200 1일 수익률 $r_t$와 원/달러 환율 수준
$z_t=(\log e_t-\mu_{t-1,252})/\sigma_{t-1,252}$이다. 평균과 표준편차는 $t-1$까지의
252거래일만 쓴다. 삼성전자의 KOSPI200 비중 때문에 종목 모멘텀과 시장수익률이
겹치며 특징--컨텍스트 상관은 $-0.29$에서 $0.22$ 범위다. 따라서 개별 계수는 다른
항이 통제된 조건부 연관이다.

\subsection{근시안적 보상과 Lasso로의 환원}
\label{sec:model}

유형 $i$의 행동점수를
\begin{equation}
q_{i,t}=(\beta_i+B_iC_t)^{\top}x_{i,t}+\alpha_i^{\top}C_t
\label{eq:score}
\end{equation}
로 둔다. $\alpha$는 종목 특징이 평균일 때도 행동 수준을 옮기는 수준효과, $B$는
컨텍스트에 따라 특징 민감도가 달라지는 기울기 효과다. 상호작용을 넣으면서 하위차수
항 $C_t$를 빼는 것은 표준적 사양 오류이므로 $\alpha$의 포함은 자료로 정당화할
대상이 아니다. 연속 행동에 대한 단기 보상을
$R_i=q_{i,t}a-\tfrac{\kappa}{2}a^{2}$로 두면 최적 정책은
\begin{equation}
a^{*}_{i,t+1}=\mathrm{clip}(q_{i,t}/\kappa,-1,1)
\label{eq:policy}
\end{equation}
이다. 이차항이 없으면 해가 거의 항상 $\pm1$이 되어 행동 강도를 재현할 수 없다.
$\kappa$는 식별을 위해 1로 고정하며, 이는 보상과 계수의 단위를 고정하는 정규화이지
곡률이 경험적으로 1이라는 주장이 아니다. 학습목적은
\begin{equation}
\min_{\theta_i}\ \frac{1}{|\mathcal{D}_i|}
\sum_{t\in\mathcal{D}_i}\!\bigl(a^{*}_{i,t+1}-a_{i,t+1}\bigr)^{2}
+\lambda_i\lVert\theta_i\rVert_1
\label{eq:loss}
\end{equation}
이고 Adam~\cite{kingma2015adam}으로 유형별 독립 학습한다.

\begin{proposition}[Lasso 환원]
\label{prop:lasso}
$z_{i,t}=[\,x_{i,t};\ \mathrm{vec}(x_{i,t}C_t^{\top});\ C_t\,]\in\mathbb{R}^{11}$,
$\theta_i=[\,\beta_i;\ \mathrm{vec}(B_i);\ \alpha_i\,]$로 두면
$q_{i,t}=\theta_i^{\top}z_{i,t}$이다. 학습 표본 전체에서 $|q_{i,t}|\le 1$이면
식~\eqref{eq:policy}의 clip이 작동하지 않아 $a^{*}=q$이고, 식~\eqref{eq:loss}은
$a_{i,t+1}$을 $z_{i,t}$에 회귀하는 $L_1$ 회귀(Lasso)와 정확히 같아진다.
\end{proposition}

본 표본에서 $|q|$의 최대값은 외국인 0.842, 개인 0.587, 기관 0.135이고 45개 분할
전부에서 포화 비율이 정확히 0이므로, 명제~\ref{prop:lasso}의 전제가 실제로
성립한다. 즉 SPR-IRL에서 ``선호를 복원한다''와 ``수급을 회귀한다''는 수치적으로
구별되지 않는다.

이 일치는 결함이 아니라 근시안 가정의 귀결이다. 한 걸음만 내다보면 오늘의 행동이
내일 이후에 영향을 주지 않으므로, 무엇을 원하는가(보상)와 무엇을 하는가(정책)가
하나로 겹친다. 최대엔트로피 IRL~\cite{ziebart2008maximum}의
$\pi\propto\exp(R/\tau)$를 적용하면 $\pi$는 평균 $q$의 절단정규분포가 되어
점추정치가 위와 일치하므로, 이 환원은 특정 구현의 우연이 아니라 근시안ㆍ이차보상
설정 전반의 성질이다. 역강화학습 정식화의 가치는 확장 경로에 있다. 상태 동역학과
다기간 가치를 넣는 순간 보상은 정책과 갈라지고 회귀로는 도달할 수 없는 구조적
복원이 된다(\S\ref{sec:discussion}).

\subsection{검증 프로토콜}
\label{sec:protocol}

복원된 계수를 그대로 보고하지 않고 네 단계를 거친다. \textbf{(V1)} 10개 시간 폴드
중 2개를 시험으로 쓰는 조합형 정화 교차검증~\cite{lopez2018advances}으로
$\binom{10}{2}=45$개 분할을 만들고, 1일 purge와 5일 embargo를 적용하며 모든
표준화는 학습 인덱스에만 적합한다. 분할들이 관측을 공유하므로 부호 일관성은
$p$-값이 아니라 안정성 지표로만 해석한다. \textbf{(V2)} 그 한계를 보완하기 위해
표본을 달력월 블록으로 나누어 1{,}000회 복원추출하고 95\% 신뢰구간을 구한다. 월
블록은 일별 수급의 자기상관을 블록 안에 보존한다. \textbf{(V3)} CPCV는 시험 폴드가
학습 폴드보다 앞설 수 있으므로, 최초 400개 관측에서 시작해 5일 purge 후 다음
60거래일을 예측하고 학습 구간을 60거래일씩 넓히는 expanding walk-forward를 10개 창
수행한다. \textbf{(V4)} 결과가 $L_1$ 벌점의 산물이 아님을 보이기 위해 $L_2$ 및
무벌점으로 재추정해 부호를 비교한다.

지표는 MAE, RMSE, 방향 정확도($\mathrm{sign}(a^{*})=\mathrm{sign}(a)$의 비율),
예측--실제 상관, 포화 비율이다. 유형 간 행동 표준편차가 다르므로 오차의 절대값은
유형 간 비교에 쓰지 않는다. 자명한 참조로 전일 자기행동을 그대로 예측하는 지속성
규칙을 함께 보고한다.

% =====================================================================
\section{실험}
\label{sec:experiments}

\subsection{데이터}
\label{sec:data}

삼성전자(005930) 일별 자료에서 252일 환율 이력, 20일 모멘텀, 다음 날 행동을 모두
계산할 수 있는 최신 973행을 사용한다. 상태일은 2022-01-06--2025-12-29이다. 단일
종목은 표본의 한계가 아니라 종목 고정효과와 유동성 차이를 배제하고 유형 간 선호
이질성만 분리하기 위한 통제된 설정이다.

행동 평균은 세 유형 모두 0에 가깝지만 표준편차는 개인 0.334, 외국인 0.257, 기관
0.125로 다르고, 1차 자기상관은 각각 0.345, 0.401, 0.149다. 동일 거래일 유형 간
행동은 강한 음의 관계를 갖는다(외국인--개인 $-0.772$, 기관--개인 $-0.499$). 이는 한
유형의 매수가 다른 유형의 매도로 나타날 수밖에 없는 시장 청산 제약이며,
\S\ref{sec:beta}의 해석 범위를 제한한다.

학습 설정은 유형별로 외국인 75 epochs / batch 256 / lr 0.001 / $\lambda=0$, 기관
10 / 2048 / 0.0005 / 0.01, 개인 20 / 512 / 0.001 / 0.0003이며 seed 42, $\kappa=1$,
유형당 파라미터 11개($\beta$ 3, $B$ 6, $\alpha$ 2)다.

\subsection{표본외 행동 재현}
\label{sec:oos}

\begin{figure}[t]
  \centering
  \includegraphics[width=\columnwidth]{figures/fig_actual_vs_prediction.pdf}
  \caption{실제 행동과 CPCV 표본외 예측(20거래일 이동평균). 각 날짜의 9개 시험예측을
  먼저 평균했다. 괄호 안은 평활 계열 간 상관으로,
  표~\ref{tab:performance}의 일별 상관과 직접 비교하지 않는다.}
  \Description{외국인, 기관, 개인의 실제 순매수와 모형 예측을 시간축에 겹쳐 그린
  세 개의 선그래프. 기관 패널에서 예측선이 0 근방에 평탄하게 놓여 있다.}
  \label{fig:prediction}
\end{figure}

\begin{table}[t]
  \caption{CPCV 표본외 성능(45개 분할 평균 $\pm$ 표준편차)과 지속성 기준선}
  \label{tab:performance}
  \small
  \centering
  \begin{tabular}{@{}lccc@{}}
    \toprule
    지표 & 외국인 & 기관 & 개인 \\
    \midrule
    \multicolumn{4}{@{}l}{\emph{SPR-IRL}}\\
    방향 정확도 & $0.636\pm0.039$ & $0.542\pm0.038$ & $0.608\pm0.044$ \\
    상관계수    & $0.278\pm0.110$ & $0.068\pm0.064$ & $0.245\pm0.118$ \\
    MAE         & $0.195\pm0.019$ & $0.094\pm0.008$ & $0.264\pm0.029$ \\
    MAE / 행동 SD & $0.759$ & $0.756$ & $0.791$ \\
    포화 비율   & $0.000$ & $0.000$ & $0.000$ \\
    \addlinespace
    \multicolumn{4}{@{}l}{\emph{지속성 기준선} ($a_{t+1}\!\leftarrow\!a_t$)}\\
    방향 정확도 & $\mathbf{0.642}$ & $\mathbf{0.547}$ & $\mathbf{0.616}$ \\
    상관계수    & $\mathbf{0.359}$ & $\mathbf{0.136}$ & $\mathbf{0.311}$ \\
    \bottomrule
  \end{tabular}
\end{table}

그림~\ref{fig:prediction}에서 외국인과 개인의 예측 계열은 실제 계열의 국면 전환을
따라가지만(평활 상관 0.69, 0.76) 진폭은 $L_1$ 축소 때문에 작다. 기관의 예측 계열은
사실상 0선에 붙어 있다(0.02).

표~\ref{tab:performance}에서 두 가지가 드러난다. 첫째, 분할 간 산포가 커서 단일
분할의 수치로 유형을 비교할 수 없다. 둘째, \textbf{자명한 지속성 기준선이 방향
정확도와 상관계수에서 SPR-IRL을 앞선다}. $a_t$는 $t$일 종가 이후 공개되므로 이는
정보상 이용 가능한 규칙이다. 원인은 명확하다. 행동의 자기상관이
$0.401/0.149/0.345$인데 보상 특징 집합에 자기 지연행동이 없다. 특징을
(F1)--(F3)로 제한한 것은 세 유형이 독립적으로 확립된 규칙성에 대응하는 특징만
공유하도록 한 설계이지 예측 성능을 위한 설계가 아니다. 표~\ref{tab:performance}은
모형이 잡음만 학습한 것이 아님을 보이는 용도이며 예측 우위의 근거가 아니다.

모든 유형ㆍ모든 분할의 포화 비율이 0이므로 명제~\ref{prop:lasso}의 전제가 성립하고,
SPR-IRL의 추정은 문자 그대로 Lasso다.

\subsection{복원된 선호: 모멘텀 축의 거울상}
\label{sec:beta}

\begin{figure*}[t]
  \centering
  \includegraphics[width=\textwidth]{figures/fig_beta_alpha_forest.pdf}
  \caption{복원된 기본 보상가중치 $\beta$(A)와 컨텍스트 주효과 $\alpha$(B)의 안정성.
  회색 점은 45개 개별 CPCV 분할, 굵은 오차선은 분할 간 표준편차, 옅은 가로선은 달력월
  블록 부트스트랩 95\% 신뢰구간이다. 속이 빈 표식은 신뢰구간이 0을 포함하는 항이다.}
  \Description{두 패널의 forest plot. 패널 A는 세 유형의 모멘텀, 군집, 손실구간
  계수를, 패널 B는 KOSPI200과 환율 주효과를 보인다. 외국인 모멘텀은 양, 개인
  모멘텀은 음이며 기관의 모든 표식은 속이 비어 있다.}
  \label{fig:forest}
\end{figure*}

\textbf{모멘텀 축의 거울상.} 외국인의 모멘텀 가중치는 $+0.0496$으로 45개 분할
전부에서 양수, 개인은 $-0.0304$로 전부에서 음수이며 두 부트스트랩 신뢰구간은 서로
겹치지 않고 각각 0을 제외한다(그림~\ref{fig:forest}A). 동일한 특징 집합을 유형별로
독립 추정했는데도 두 유형이 정반대로 갈리며, 이는 선행
연구~\cite{grinblatt2000investment,choe1999foreign}의 부호 기대와 일치한다.

\textbf{컨텍스트에서의 재현.} 같은 거울상이 주효과에서도 나타난다
(그림~\ref{fig:forest}B). 외국인은 시장 상승일에 순매수 방향으로($+0.0361$), 원화
약세 국면에 순매도 방향으로($-0.0261$) 이동하고 개인은 두 컨텍스트 모두에서 정확히
반대다($-0.0246$, $+0.0298$). 네 계수 모두 45개 분할 전부에서 부호가 유지되고
신뢰구간도 0을 제외한다.

\textbf{군집.} 군집 계수는 세 유형 모두 45개 분할 전부에서 음수이고 외국인
($-0.0402$)과 개인($-0.0322$)은 신뢰구간도 0을 제외한다. 그러나 이 결과는 행동적
해석과 시장 청산 제약을 분리하지 못한다. \S\ref{sec:data}의 $-0.77$이라는 동시점
상관 아래에서는 음의 군집 계수 중 상당 부분이 회계적 제약에서 발생할 수 있으므로,
본 연구는 이를 ``관측된 집계 수급이 유형 간 음의 교차 연관을 보인다''는 서술적
결과로 제한한다. 거울상 주효과에도 같은 유보가 적용된다.

\textbf{손실구간.} 개인의 손실구간 가중치는 $+0.0293$으로 45개 분할 전부에서
양수이며 처분효과 문헌~\cite{shefrin1985disposition,odean1998investors}의 기준가격
의존성과 정합한다. 그러나 그림~\ref{fig:forest}A에서 보듯 세 유형 모두 신뢰구간이
0을 포함한다. 개인의 구간은 $[-0.0122,+0.0665]$로, 45개 분할에서 부호가 100\%
일관됨에도 월 블록 재표집에서는 5.5\%가 부호를 뒤집는다. 이는 CPCV 분할이 관측을
공유하기 때문에 부호 일관성이 불확실성을 과소평가한다는 \S\ref{sec:protocol}의
경고가 실제로 작동한 사례다. 따라서 처분효과와 정합하는 부호를 관측했다고 보고하되
통계적으로 확립된 결과로 주장하지 않는다.

상호작용 $B$ 12개 중 45개 분할 전부에서 부호가 유지된 것은 외국인
모멘텀$\times$KOSPI200($-0.0139$), 개인 모멘텀$\times$환율($+0.0144$), 개인
군집$\times$KOSPI200($-0.0080$), 개인 손실구간$\times$환율($-0.0186$) 넷이고, 이 중
부트스트랩을 통과하는 것은 첫째와 넷째다. 12개를 동시에 검토했으므로 상호작용은
탐색적 결과로만 제시한다.

\subsection{안정성}
\label{sec:stability}

\begin{figure*}[t]
  \centering
  \includegraphics[width=\textwidth]{figures/fig_walk_forward.pdf}
  \caption{Expanding walk-forward 10개 창의 계수 경로와 방향 정확도. 성능은 창마다
  흔들리지만 계수 부호는 유지되며, 외국인과 개인의 모멘텀 경로는 한 번도 교차하지
  않는다.}
  \Description{확장 창 번호를 가로축으로 하여 세 유형의 모멘텀 계수, 두 컨텍스트
  주효과, 방향 정확도를 그린 네 개의 선그래프.}
  \label{fig:walkforward}
\end{figure*}

\textbf{(V2)} 부트스트랩 결과는 셋으로 갈린다. 모멘텀 거울상과 컨텍스트 거울상을
이루는 8개 항이 모두 0을 제외하고, 손실구간은 통과하지 못하며, 기관은 어떤 항도
0을 제외하지 못한다(그림~\ref{fig:forest}의 속 빈 표식).

\textbf{(V3)} 그림~\ref{fig:walkforward}에서 성능은 크게 흔들린다. 외국인 방향
정확도는 0.43--0.85를 오가며 평균 $0.610\pm0.116$, 상관은 $0.131\pm0.162$로 CPCV의
0.278에서 절반 이하다. CPCV는 시험 폴드를 학습 폴드보다 앞세울 수 있는 반면
walk-forward는 그렇지 않으므로, 본 연구는 이 저하를 시간 구조에서 오는 일반화
한계로 그대로 보고한다. 반면 계수 경로는 흔들리지 않는다. 외국인 모멘텀은 10개 창
전부 양수, 개인은 전부 음수이며 두 경로가 한 번도 교차하지 않고, 두 컨텍스트
주효과에서도 같은 거울상이 유지된다. 기관의 경로는 0선에 붙어 있다.

\textbf{(V4)} 외국인과 개인의 계수는 $L_1$, $L_2$, 무벌점에서 소수점 넷째 자리까지
일치하므로 거울상은 벌점의 산물이 아니다. 반면 기관의 계수만 벌점에 따라
이동한다(군집 $-0.0076\rightarrow-0.0124$, 손실구간
$-0.0000\rightarrow-0.0018$).

\subsection{기관: 수렴된 null}
\label{sec:null}

기관 수급에 대해 SPR-IRL은 일별 빈도에서 안정적으로 식별되는 선호를 찾지 못했고,
본 연구는 이를 은폐하지 않고 결과로 보고한다. 근거는 다섯이다. (i) 예측 계열이 0선에
퇴화한다(그림~\ref{fig:prediction}). (ii) 군집을 제외한 계수가 다른 유형의 $1/40$
수준이다($-0.0006$, $-0.0008$). (iii) 부트스트랩에서 0을 제외하는 항이
없다(그림~\ref{fig:forest}). (iv) 계수가 벌점 선택에 좌우된다. (v) walk-forward
10개 창에서 모멘텀 부호 일관성이 40\%, 손실구간이 0\%다.

\textbf{이 null은 최적화 실패가 아니다.} 명제~\ref{prop:lasso}에 따라
식~\eqref{eq:loss}은 볼록 문제이므로 전역 최적해가 유일하다. 도달한 해와 동일
목적함수의 폐쇄형 최적해를 비교하면 기관의 목적함수 상대 격차는 평균 0.17\%, 최대
0.48\%에 불과하다(외국인 0.21\%, 개인 0.92\%). 세 유형 모두 도달 가능한 최적해에
근접했고, 기관에서만 그 최적해의 계수가 0 근방이다. 데이터가 기관에 대해 결정하는
선호가 없는 것이지 최적화가 찾지 못한 것이 아니다.

경제적으로 이 null은 자연스럽다. ``기관''은 연기금, 투신, 보험, 사모, 금융투자를
합산한 범주이므로 서로 다른 목적함수를 가진 주체들의 수급이 상쇄된다.

\subsection{구성 타당성 점검}
\label{sec:construct}

\begin{table}[t]
  \caption{사전 부호 기대와 복원 결과}
  \label{tab:construct}
  \small
  \begin{tabular}{@{}>{\raggedright\arraybackslash}p{0.30\columnwidth}
                  >{\raggedright\arraybackslash}p{0.13\columnwidth}
                  >{\raggedright\arraybackslash}p{0.26\columnwidth}
                  >{\raggedright\arraybackslash}p{0.16\columnwidth}@{}}
    \toprule
    규칙성 & 사전 기대 & 복원 결과 & 판정 \\
    \midrule
    외국인 추세거래~\cite{grinblatt2000investment,choe1999foreign}
      & $\beta^{mom}\!>\!0$ & $+0.0496$ (100\%, CI 0 제외) & 부합 \\
    개인 역추세~\cite{grinblatt2000investment,kaniel2008individual}
      & $\beta^{mom}\!<\!0$ & $-0.0304$ (100\%, CI 0 제외) & 부합 \\
    처분효과~\cite{shefrin1985disposition,odean1998investors}
      & $\beta^{uw}\!>\!0$ & $+0.0293$ (100\%, CI 0 포함) & 부분 부합 \\
    기관 군집행동~\cite{lakonishok1992impact}
      & $\beta^{herd}\!>\!0$ & $-0.0044$ (100\%, CI 0 포함) & 불일치 \\
    \bottomrule
  \end{tabular}
\end{table}

표~\ref{tab:construct}이 대조 결과다. 동일한 특징과 모형 형태를 부여하고 파라미터만
분리 학습했음에도 복원된 부호가 독립적으로 확립된 규칙성과 부합하고 두 유형이
모멘텀에서 정반대로 갈린다는 점이 구성 타당성의 증거다. 불일치 항목은 측정 대상이
다르기 때문이다. Lakonishok 외~\cite{lakonishok1992impact}는 기관 \emph{내부}의
군집을, 식~\eqref{eq:herd}은 기관이 \emph{다른 두 유형}과 같은 방향으로 움직이는지를
측정하므로 본 결과는 해당 문헌의 반증이 아니다.

%  패키지: kotex, booktabs, graphicx, amsmath, array
%  그림:  figures/  (make_paper_figures.py 산출물 중 3개만 사용)
%  참고문헌: paper-ko.bib + paper-refs-additional.bib
%
%  분량 설계: 3장 약 1.6쪽 + 4장 약 2.4쪽 = 약 4쪽.
%  1~2장에 3쪽, 5장+참고문헌에 1쪽을 남긴다.
% =====================================================================

\section{방법: SPR-IRL}
\label{sec:method}

본 절은 대칭 선호복원(Symmetric Preference Recovery, 이하 SPR-IRL) 절차를
정의한다. SPR-IRL은 모형 하나가 아니라 (i) 공유 특징 위의 대칭 추정, (ii) 표본외
안정성 평가, (iii) 사전 부호 기대와의 대조로 이루어진 절차다.

\subsection{연속 행동과 대칭 추정}
\label{sec:setup}

투자자 유형을 $\mathcal{I}=\{\text{외국인},\text{기관},\text{개인}\}$으로 둔다.
유형 $i$의 $t$일 총매수ㆍ총매도금액을 $V^{buy}_{i,t}$, $V^{sell}_{i,t}$라 하면
행동을
\begin{equation}
u_{i,t}=\frac{V^{buy}_{i,t}-V^{sell}_{i,t}}
             {V^{buy}_{i,t}+V^{sell}_{i,t}}\in[-1,1]
\label{eq:action}
\end{equation}
로 정의한다. 분모가 시장 전체가 아니라 해당 유형의 총거래금액이므로 $u_{i,t}$는
유형별 거래 규모에 영향받지 않는 공통 척도가 되며, 이진 라벨과 달리 강도 정보를
보존한다.

세 유형은 동일한 상태 특징 $x_{i,t}\in\mathbb{R}^{3}$, 동일한 컨텍스트
$C_t\in\mathbb{R}^{2}$, 동일한 모형 형태를 부여받고 파라미터
$\theta_i=(\beta_i,B_i,\alpha_i)$만 각자의 데이터로 분리 학습한다. 유형별로 다른
특징을 미리 배정하지 않으므로, 결과의 유형 간 차이는 설계가 아니라 데이터에서
학습된 것이며 이것이 계수의 직접 비교를 정당화한다.

학습 표본은 $(x_{i,t},C_t,a_{i,t+1})$, $a_{i,t+1}=u_{i,t+1}$이다. $t$일 상태가
$t+1$일 행동을 설명하므로 미래 거래는 특징에 들어가지 않는다.

\subsection{보상 특징과 시장 컨텍스트}
\label{sec:features}

세 특징과 두 컨텍스트는 모두 각 분할의 학습 인덱스에서만 표준화하므로, 이후 계수는
1표준편차 변화의 효과로 읽는다.

\textbf{(F1) 중기 모멘텀} $x^{mom}_{t}=\log(P_t/P_{t-20})$. 20거래일은 해석을 위한
사전 지정값이며 자료에서 최적화하지 않았다. 양수는 추세추종, 음수는 역추세를 뜻하고
부호 기대는 유형별 모멘텀 거래 문헌~\cite{grinblatt2000investment,choe1999foreign}에서
온다.

\textbf{(F2) 군집행동} 유형 $i$를 제외한 두 유형의 전일 행동 횡단면 평균
\begin{equation}
x^{herd}_{i,t}=\tfrac{1}{2}\left(u_{j,t-1}+u_{k,t-1}\right),\qquad j,k\neq i.
\label{eq:herd}
\end{equation}
1거래일 지연은 정보 제약이 아니라, 동일 거래일 유형 간 수급의 회계적 상쇄를 예측
변수에서 배제하기 위한 설계 선택이다(\S\ref{sec:beta}). 부호 기대는 기관 군집행동
문헌~\cite{lakonishok1992impact}에서 오나, 해당 문헌은 기관 \emph{내부}의 군집을,
식~\eqref{eq:herd}은 \emph{유형 간} 관계를 측정한다.

\textbf{(F3) 평균단가 대비 손실구간} 평균단가 대용치 $\bar P_{i,t}$와 당일 총거래
VWAP $P^{vwap}_{i,t}$로
\begin{equation}
x^{uw}_{i,t}=\max\!\left(
\frac{\bar P_{i,t}-P^{vwap}_{i,t}}{\bar P_{i,t}},\,0\right).
\label{eq:underwater}
\end{equation}
0에서 절단하므로 이 변수는 준거점 의존적 설계다. 평균단가는 총매수ㆍ총매도 수량과
금액으로 유효 보유량을 재귀 갱신해 구하되 전일 보유량에 망각계수 $\rho=0.98$을
곱한다(반감기 약 34거래일). 부호 기대는 처분효과
문헌~\cite{shefrin1985disposition,odean1998investors}에서 온다. 다만 이 대용치는
집계 거래로 만든 근사치이므로 $x^{uw}$의 계수를 처분효과의 직접 검정으로 읽을 수
없다.

\textbf{컨텍스트}는 KOSPI200 1일 수익률 $r_t$와 원/달러 환율 수준
$z_t=(\log e_t-\mu_{t-1,252})/\sigma_{t-1,252}$이다. 평균과 표준편차는 $t-1$까지의
252거래일만 쓴다. 삼성전자의 KOSPI200 비중 때문에 종목 모멘텀과 시장수익률이
겹치며 특징--컨텍스트 상관은 $-0.29$에서 $0.22$ 범위다. 따라서 개별 계수는 다른
항이 통제된 조건부 연관이다.

\subsection{근시안적 보상과 Lasso로의 환원}
\label{sec:model}

유형 $i$의 행동점수를
\begin{equation}
q_{i,t}=(\beta_i+B_iC_t)^{\top}x_{i,t}+\alpha_i^{\top}C_t
\label{eq:score}
\end{equation}
로 둔다. $\alpha$는 종목 특징이 평균일 때도 행동 수준을 옮기는 수준효과, $B$는
컨텍스트에 따라 특징 민감도가 달라지는 기울기 효과다. 상호작용을 넣으면서 하위차수
항 $C_t$를 빼는 것은 표준적 사양 오류이므로 $\alpha$의 포함은 자료로 정당화할
대상이 아니다. 연속 행동에 대한 단기 보상을
$R_i=q_{i,t}a-\tfrac{\kappa}{2}a^{2}$로 두면 최적 정책은
\begin{equation}
a^{*}_{i,t+1}=\mathrm{clip}(q_{i,t}/\kappa,-1,1)
\label{eq:policy}
\end{equation}
이다. 이차항이 없으면 해가 거의 항상 $\pm1$이 되어 행동 강도를 재현할 수 없다.
$\kappa$는 식별을 위해 1로 고정하며, 이는 보상과 계수의 단위를 고정하는 정규화이지
곡률이 경험적으로 1이라는 주장이 아니다. 학습목적은
\begin{equation}
\min_{\theta_i}\ \frac{1}{|\mathcal{D}_i|}
\sum_{t\in\mathcal{D}_i}\!\bigl(a^{*}_{i,t+1}-a_{i,t+1}\bigr)^{2}
+\lambda_i\lVert\theta_i\rVert_1
\label{eq:loss}
\end{equation}
이고 Adam~\cite{kingma2015adam}으로 유형별 독립 학습한다.

\begin{proposition}[Lasso 환원]
\label{prop:lasso}
$z_{i,t}=[\,x_{i,t};\ \mathrm{vec}(x_{i,t}C_t^{\top});\ C_t\,]\in\mathbb{R}^{11}$,
$\theta_i=[\,\beta_i;\ \mathrm{vec}(B_i);\ \alpha_i\,]$로 두면
$q_{i,t}=\theta_i^{\top}z_{i,t}$이다. 학습 표본 전체에서 $|q_{i,t}|\le 1$이면
식~\eqref{eq:policy}의 clip이 작동하지 않아 $a^{*}=q$이고, 식~\eqref{eq:loss}은
$a_{i,t+1}$을 $z_{i,t}$에 회귀하는 $L_1$ 회귀(Lasso)와 정확히 같아진다.
\end{proposition}

본 표본에서 $|q|$의 최대값은 외국인 0.842, 개인 0.587, 기관 0.135이고 45개 분할
전부에서 포화 비율이 정확히 0이므로, 명제~\ref{prop:lasso}의 전제가 실제로
성립한다. 즉 SPR-IRL에서 ``선호를 복원한다''와 ``수급을 회귀한다''는 수치적으로
구별되지 않는다.

이 일치는 결함이 아니라 근시안 가정의 귀결이다. 한 걸음만 내다보면 오늘의 행동이
내일 이후에 영향을 주지 않으므로, 무엇을 원하는가(보상)와 무엇을 하는가(정책)가
하나로 겹친다. 최대엔트로피 IRL~\cite{ziebart2008maximum}의
$\pi\propto\exp(R/\tau)$를 적용하면 $\pi$는 평균 $q$의 절단정규분포가 되어
점추정치가 위와 일치하므로, 이 환원은 특정 구현의 우연이 아니라 근시안ㆍ이차보상
설정 전반의 성질이다. 역강화학습 정식화의 가치는 확장 경로에 있다. 상태 동역학과
다기간 가치를 넣는 순간 보상은 정책과 갈라지고 회귀로는 도달할 수 없는 구조적
복원이 된다(\S\ref{sec:discussion}).

\subsection{검증 프로토콜}
\label{sec:protocol}

복원된 계수를 그대로 보고하지 않고 네 단계를 거친다. \textbf{(V1)} 10개 시간 폴드
중 2개를 시험으로 쓰는 조합형 정화 교차검증~\cite{lopez2018advances}으로
$\binom{10}{2}=45$개 분할을 만들고, 1일 purge와 5일 embargo를 적용하며 모든
표준화는 학습 인덱스에만 적합한다. 분할들이 관측을 공유하므로 부호 일관성은
$p$-값이 아니라 안정성 지표로만 해석한다. \textbf{(V2)} 그 한계를 보완하기 위해
표본을 달력월 블록으로 나누어 1{,}000회 복원추출하고 95\% 신뢰구간을 구한다. 월
블록은 일별 수급의 자기상관을 블록 안에 보존한다. \textbf{(V3)} CPCV는 시험 폴드가
학습 폴드보다 앞설 수 있으므로, 최초 400개 관측에서 시작해 5일 purge 후 다음
60거래일을 예측하고 학습 구간을 60거래일씩 넓히는 expanding walk-forward를 10개 창
수행한다. \textbf{(V4)} 결과가 $L_1$ 벌점의 산물이 아님을 보이기 위해 $L_2$ 및
무벌점으로 재추정해 부호를 비교한다.

지표는 MAE, RMSE, 방향 정확도($\mathrm{sign}(a^{*})=\mathrm{sign}(a)$의 비율),
예측--실제 상관, 포화 비율이다. 유형 간 행동 표준편차가 다르므로 오차의 절대값은
유형 간 비교에 쓰지 않는다. 자명한 참조로 전일 자기행동을 그대로 예측하는 지속성
규칙을 함께 보고한다.

% =====================================================================
\section{실험}
\label{sec:experiments}

\subsection{데이터}
\label{sec:data}

삼성전자(005930) 일별 자료에서 252일 환율 이력, 20일 모멘텀, 다음 날 행동을 모두
계산할 수 있는 최신 973행을 사용한다. 상태일은 2022-01-06--2025-12-29이다. 단일
종목은 표본의 한계가 아니라 종목 고정효과와 유동성 차이를 배제하고 유형 간 선호
이질성만 분리하기 위한 통제된 설정이다.

행동 평균은 세 유형 모두 0에 가깝지만 표준편차는 개인 0.334, 외국인 0.257, 기관
0.125로 다르고, 1차 자기상관은 각각 0.345, 0.401, 0.149다. 동일 거래일 유형 간
행동은 강한 음의 관계를 갖는다(외국인--개인 $-0.772$, 기관--개인 $-0.499$). 이는 한
유형의 매수가 다른 유형의 매도로 나타날 수밖에 없는 시장 청산 제약이며,
\S\ref{sec:beta}의 해석 범위를 제한한다.

학습 설정은 유형별로 외국인 75 epochs / batch 256 / lr 0.001 / $\lambda=0$, 기관
10 / 2048 / 0.0005 / 0.01, 개인 20 / 512 / 0.001 / 0.0003이며 seed 42, $\kappa=1$,
유형당 파라미터 11개($\beta$ 3, $B$ 6, $\alpha$ 2)다.

\subsection{표본외 행동 재현}
\label{sec:oos}

\begin{figure}[t]
  \centering
  \includegraphics[width=\columnwidth]{figures/fig_actual_vs_prediction.pdf}
  \caption{실제 행동과 CPCV 표본외 예측(20거래일 이동평균). 각 날짜의 9개 시험예측을
  먼저 평균했다. 괄호 안은 평활 계열 간 상관으로,
  표~\ref{tab:performance}의 일별 상관과 직접 비교하지 않는다.}
  \Description{외국인, 기관, 개인의 실제 순매수와 모형 예측을 시간축에 겹쳐 그린
  세 개의 선그래프. 기관 패널에서 예측선이 0 근방에 평탄하게 놓여 있다.}
  \label{fig:prediction}
\end{figure}

\begin{table}[t]
  \caption{CPCV 표본외 성능(45개 분할 평균 $\pm$ 표준편차)과 지속성 기준선}
  \label{tab:performance}
  \small
  \centering
  \begin{tabular}{@{}lccc@{}}
    \toprule
    지표 & 외국인 & 기관 & 개인 \\
    \midrule
    \multicolumn{4}{@{}l}{\emph{SPR-IRL}}\\
    방향 정확도 & $0.636\pm0.039$ & $0.542\pm0.038$ & $0.608\pm0.044$ \\
    상관계수    & $0.278\pm0.110$ & $0.068\pm0.064$ & $0.245\pm0.118$ \\
    MAE         & $0.195\pm0.019$ & $0.094\pm0.008$ & $0.264\pm0.029$ \\
    MAE / 행동 SD & $0.759$ & $0.756$ & $0.791$ \\
    포화 비율   & $0.000$ & $0.000$ & $0.000$ \\
    \addlinespace
    \multicolumn{4}{@{}l}{\emph{지속성 기준선} ($a_{t+1}\!\leftarrow\!a_t$)}\\
    방향 정확도 & $\mathbf{0.642}$ & $\mathbf{0.547}$ & $\mathbf{0.616}$ \\
    상관계수    & $\mathbf{0.359}$ & $\mathbf{0.136}$ & $\mathbf{0.311}$ \\
    \bottomrule
  \end{tabular}
\end{table}

그림~\ref{fig:prediction}에서 외국인과 개인의 예측 계열은 실제 계열의 국면 전환을
따라가지만(평활 상관 0.69, 0.76) 진폭은 $L_1$ 축소 때문에 작다. 기관의 예측 계열은
사실상 0선에 붙어 있다(0.02).

표~\ref{tab:performance}에서 두 가지가 드러난다. 첫째, 분할 간 산포가 커서 단일
분할의 수치로 유형을 비교할 수 없다. 둘째, \textbf{자명한 지속성 기준선이 방향
정확도와 상관계수에서 SPR-IRL을 앞선다}. $a_t$는 $t$일 종가 이후 공개되므로 이는
정보상 이용 가능한 규칙이다. 원인은 명확하다. 행동의 자기상관이
$0.401/0.149/0.345$인데 보상 특징 집합에 자기 지연행동이 없다. 특징을
(F1)--(F3)로 제한한 것은 세 유형이 독립적으로 확립된 규칙성에 대응하는 특징만
공유하도록 한 설계이지 예측 성능을 위한 설계가 아니다. 표~\ref{tab:performance}은
모형이 잡음만 학습한 것이 아님을 보이는 용도이며 예측 우위의 근거가 아니다.

모든 유형ㆍ모든 분할의 포화 비율이 0이므로 명제~\ref{prop:lasso}의 전제가 성립하고,
SPR-IRL의 추정은 문자 그대로 Lasso다.

\subsection{복원된 선호: 모멘텀 축의 거울상}
\label{sec:beta}

\begin{figure*}[t]
  \centering
  \includegraphics[width=\textwidth]{figures/fig_beta_alpha_forest.pdf}
  \caption{복원된 기본 보상가중치 $\beta$(A)와 컨텍스트 주효과 $\alpha$(B)의 안정성.
  회색 점은 45개 개별 CPCV 분할, 굵은 오차선은 분할 간 표준편차, 옅은 가로선은 달력월
  블록 부트스트랩 95\% 신뢰구간이다. 속이 빈 표식은 신뢰구간이 0을 포함하는 항이다.}
  \Description{두 패널의 forest plot. 패널 A는 세 유형의 모멘텀, 군집, 손실구간
  계수를, 패널 B는 KOSPI200과 환율 주효과를 보인다. 외국인 모멘텀은 양, 개인
  모멘텀은 음이며 기관의 모든 표식은 속이 비어 있다.}
  \label{fig:forest}
\end{figure*}

\textbf{모멘텀 축의 거울상.} 외국인의 모멘텀 가중치는 $+0.0496$으로 45개 분할
전부에서 양수, 개인은 $-0.0304$로 전부에서 음수이며 두 부트스트랩 신뢰구간은 서로
겹치지 않고 각각 0을 제외한다(그림~\ref{fig:forest}A). 동일한 특징 집합을 유형별로
독립 추정했는데도 두 유형이 정반대로 갈리며, 이는 선행
연구~\cite{grinblatt2000investment,choe1999foreign}의 부호 기대와 일치한다.

\textbf{컨텍스트에서의 재현.} 같은 거울상이 주효과에서도 나타난다
(그림~\ref{fig:forest}B). 외국인은 시장 상승일에 순매수 방향으로($+0.0361$), 원화
약세 국면에 순매도 방향으로($-0.0261$) 이동하고 개인은 두 컨텍스트 모두에서 정확히
반대다($-0.0246$, $+0.0298$). 네 계수 모두 45개 분할 전부에서 부호가 유지되고
신뢰구간도 0을 제외한다.

\textbf{군집.} 군집 계수는 세 유형 모두 45개 분할 전부에서 음수이고 외국인
($-0.0402$)과 개인($-0.0322$)은 신뢰구간도 0을 제외한다. 그러나 이 결과는 행동적
해석과 시장 청산 제약을 분리하지 못한다. \S\ref{sec:data}의 $-0.77$이라는 동시점
상관 아래에서는 음의 군집 계수 중 상당 부분이 회계적 제약에서 발생할 수 있으므로,
본 연구는 이를 ``관측된 집계 수급이 유형 간 음의 교차 연관을 보인다''는 서술적
결과로 제한한다. 거울상 주효과에도 같은 유보가 적용된다.

\textbf{손실구간.} 개인의 손실구간 가중치는 $+0.0293$으로 45개 분할 전부에서
양수이며 처분효과 문헌~\cite{shefrin1985disposition,odean1998investors}의 기준가격
의존성과 정합한다. 그러나 그림~\ref{fig:forest}A에서 보듯 세 유형 모두 신뢰구간이
0을 포함한다. 개인의 구간은 $[-0.0122,+0.0665]$로, 45개 분할에서 부호가 100\%
일관됨에도 월 블록 재표집에서는 5.5\%가 부호를 뒤집는다. 이는 CPCV 분할이 관측을
공유하기 때문에 부호 일관성이 불확실성을 과소평가한다는 \S\ref{sec:protocol}의
경고가 실제로 작동한 사례다. 따라서 처분효과와 정합하는 부호를 관측했다고 보고하되
통계적으로 확립된 결과로 주장하지 않는다.

상호작용 $B$ 12개 중 45개 분할 전부에서 부호가 유지된 것은 외국인
모멘텀$\times$KOSPI200($-0.0139$), 개인 모멘텀$\times$환율($+0.0144$), 개인
군집$\times$KOSPI200($-0.0080$), 개인 손실구간$\times$환율($-0.0186$) 넷이고, 이 중
부트스트랩을 통과하는 것은 첫째와 넷째다. 12개를 동시에 검토했으므로 상호작용은
탐색적 결과로만 제시한다.

\subsection{안정성}
\label{sec:stability}

\begin{figure*}[t]
  \centering
  \includegraphics[width=\textwidth]{figures/fig_walk_forward.pdf}
  \caption{Expanding walk-forward 10개 창의 계수 경로와 방향 정확도. 성능은 창마다
  흔들리지만 계수 부호는 유지되며, 외국인과 개인의 모멘텀 경로는 한 번도 교차하지
  않는다.}
  \Description{확장 창 번호를 가로축으로 하여 세 유형의 모멘텀 계수, 두 컨텍스트
  주효과, 방향 정확도를 그린 네 개의 선그래프.}
  \label{fig:walkforward}
\end{figure*}

\textbf{(V2)} 부트스트랩 결과는 셋으로 갈린다. 모멘텀 거울상과 컨텍스트 거울상을
이루는 8개 항이 모두 0을 제외하고, 손실구간은 통과하지 못하며, 기관은 어떤 항도
0을 제외하지 못한다(그림~\ref{fig:forest}의 속 빈 표식).

\textbf{(V3)} 그림~\ref{fig:walkforward}에서 성능은 크게 흔들린다. 외국인 방향
정확도는 0.43--0.85를 오가며 평균 $0.610\pm0.116$, 상관은 $0.131\pm0.162$로 CPCV의
0.278에서 절반 이하다. CPCV는 시험 폴드를 학습 폴드보다 앞세울 수 있는 반면
walk-forward는 그렇지 않으므로, 본 연구는 이 저하를 시간 구조에서 오는 일반화
한계로 그대로 보고한다. 반면 계수 경로는 흔들리지 않는다. 외국인 모멘텀은 10개 창
전부 양수, 개인은 전부 음수이며 두 경로가 한 번도 교차하지 않고, 두 컨텍스트
주효과에서도 같은 거울상이 유지된다. 기관의 경로는 0선에 붙어 있다.

\textbf{(V4)} 외국인과 개인의 계수는 $L_1$, $L_2$, 무벌점에서 소수점 넷째 자리까지
일치하므로 거울상은 벌점의 산물이 아니다. 반면 기관의 계수만 벌점에 따라
이동한다(군집 $-0.0076\rightarrow-0.0124$, 손실구간
$-0.0000\rightarrow-0.0018$).

\subsection{기관: 수렴된 null}
\label{sec:null}

기관 수급에 대해 SPR-IRL은 일별 빈도에서 안정적으로 식별되는 선호를 찾지 못했고,
본 연구는 이를 은폐하지 않고 결과로 보고한다. 근거는 다섯이다. (i) 예측 계열이 0선에
퇴화한다(그림~\ref{fig:prediction}). (ii) 군집을 제외한 계수가 다른 유형의 $1/40$
수준이다($-0.0006$, $-0.0008$). (iii) 부트스트랩에서 0을 제외하는 항이
없다(그림~\ref{fig:forest}). (iv) 계수가 벌점 선택에 좌우된다. (v) walk-forward
10개 창에서 모멘텀 부호 일관성이 40\%, 손실구간이 0\%다.

\textbf{이 null은 최적화 실패가 아니다.} 명제~\ref{prop:lasso}에 따라
식~\eqref{eq:loss}은 볼록 문제이므로 전역 최적해가 유일하다. 도달한 해와 동일
목적함수의 폐쇄형 최적해를 비교하면 기관의 목적함수 상대 격차는 평균 0.17\%, 최대
0.48\%에 불과하다(외국인 0.21\%, 개인 0.92\%). 세 유형 모두 도달 가능한 최적해에
근접했고, 기관에서만 그 최적해의 계수가 0 근방이다. 데이터가 기관에 대해 결정하는
선호가 없는 것이지 최적화가 찾지 못한 것이 아니다.

경제적으로 이 null은 자연스럽다. ``기관''은 연기금, 투신, 보험, 사모, 금융투자를
합산한 범주이므로 서로 다른 목적함수를 가진 주체들의 수급이 상쇄된다.

\subsection{구성 타당성 점검}
\label{sec:construct}

\begin{table}[t]
  \caption{사전 부호 기대와 복원 결과}
  \label{tab:construct}
  \small
  \begin{tabular}{@{}>{\raggedright\arraybackslash}p{0.30\columnwidth}
                  >{\raggedright\arraybackslash}p{0.13\columnwidth}
                  >{\raggedright\arraybackslash}p{0.26\columnwidth}
                  >{\raggedright\arraybackslash}p{0.16\columnwidth}@{}}
    \toprule
    규칙성 & 사전 기대 & 복원 결과 & 판정 \\
    \midrule
    외국인 추세거래~\cite{grinblatt2000investment,choe1999foreign}
      & $\beta^{mom}\!>\!0$ & $+0.0496$ (100\%, CI 0 제외) & 부합 \\
    개인 역추세~\cite{grinblatt2000investment,kaniel2008individual}
      & $\beta^{mom}\!<\!0$ & $-0.0304$ (100\%, CI 0 제외) & 부합 \\
    처분효과~\cite{shefrin1985disposition,odean1998investors}
      & $\beta^{uw}\!>\!0$ & $+0.0293$ (100\%, CI 0 포함) & 부분 부합 \\
    기관 군집행동~\cite{lakonishok1992impact}
      & $\beta^{herd}\!>\!0$ & $-0.0044$ (100\%, CI 0 포함) & 불일치 \\
    \bottomrule
  \end{tabular}
\end{table}

표~\ref{tab:construct}이 대조 결과다. 동일한 특징과 모형 형태를 부여하고 파라미터만
분리 학습했음에도 복원된 부호가 독립적으로 확립된 규칙성과 부합하고 두 유형이
모멘텀에서 정반대로 갈린다는 점이 구성 타당성의 증거다. 불일치 항목은 측정 대상이
다르기 때문이다. Lakonishok 외~\cite{lakonishok1992impact}는 기관 \emph{내부}의
군집을, 식~\eqref{eq:herd}은 기관이 \emph{다른 두 유형}과 같은 방향으로 움직이는지를
측정하므로 본 결과는 해당 문헌의 반증이 아니다.

%  패키지: kotex, booktabs, graphicx, amsmath, array
%  그림:  figures/  (make_paper_figures.py 산출물 중 3개만 사용)
%  참고문헌: paper-ko.bib + paper-refs-additional.bib
%
%  분량 설계: 3장 약 1.6쪽 + 4장 약 2.4쪽 = 약 4쪽.
%  1~2장에 3쪽, 5장+참고문헌에 1쪽을 남긴다.
% =====================================================================

\section{방법: SPR-IRL}
\label{sec:method}

본 절은 대칭 선호복원(Symmetric Preference Recovery, 이하 SPR-IRL) 절차를
정의한다. SPR-IRL은 모형 하나가 아니라 (i) 공유 특징 위의 대칭 추정, (ii) 표본외
안정성 평가, (iii) 사전 부호 기대와의 대조로 이루어진 절차다.

\subsection{연속 행동과 대칭 추정}
\label{sec:setup}

투자자 유형을 $\mathcal{I}=\{\text{외국인},\text{기관},\text{개인}\}$으로 둔다.
유형 $i$의 $t$일 총매수ㆍ총매도금액을 $V^{buy}_{i,t}$, $V^{sell}_{i,t}$라 하면
행동을
\begin{equation}
u_{i,t}=\frac{V^{buy}_{i,t}-V^{sell}_{i,t}}
             {V^{buy}_{i,t}+V^{sell}_{i,t}}\in[-1,1]
\label{eq:action}
\end{equation}
로 정의한다. 분모가 시장 전체가 아니라 해당 유형의 총거래금액이므로 $u_{i,t}$는
유형별 거래 규모에 영향받지 않는 공통 척도가 되며, 이진 라벨과 달리 강도 정보를
보존한다.

세 유형은 동일한 상태 특징 $x_{i,t}\in\mathbb{R}^{3}$, 동일한 컨텍스트
$C_t\in\mathbb{R}^{2}$, 동일한 모형 형태를 부여받고 파라미터
$\theta_i=(\beta_i,B_i,\alpha_i)$만 각자의 데이터로 분리 학습한다. 유형별로 다른
특징을 미리 배정하지 않으므로, 결과의 유형 간 차이는 설계가 아니라 데이터에서
학습된 것이며 이것이 계수의 직접 비교를 정당화한다.

학습 표본은 $(x_{i,t},C_t,a_{i,t+1})$, $a_{i,t+1}=u_{i,t+1}$이다. $t$일 상태가
$t+1$일 행동을 설명하므로 미래 거래는 특징에 들어가지 않는다.

\subsection{보상 특징과 시장 컨텍스트}
\label{sec:features}

세 특징과 두 컨텍스트는 모두 각 분할의 학습 인덱스에서만 표준화하므로, 이후 계수는
1표준편차 변화의 효과로 읽는다.

\textbf{(F1) 중기 모멘텀} $x^{mom}_{t}=\log(P_t/P_{t-20})$. 20거래일은 해석을 위한
사전 지정값이며 자료에서 최적화하지 않았다. 양수는 추세추종, 음수는 역추세를 뜻하고
부호 기대는 유형별 모멘텀 거래 문헌~\cite{grinblatt2000investment,choe1999foreign}에서
온다.

\textbf{(F2) 군집행동} 유형 $i$를 제외한 두 유형의 전일 행동 횡단면 평균
\begin{equation}
x^{herd}_{i,t}=\tfrac{1}{2}\left(u_{j,t-1}+u_{k,t-1}\right),\qquad j,k\neq i.
\label{eq:herd}
\end{equation}
1거래일 지연은 정보 제약이 아니라, 동일 거래일 유형 간 수급의 회계적 상쇄를 예측
변수에서 배제하기 위한 설계 선택이다(\S\ref{sec:beta}). 부호 기대는 기관 군집행동
문헌~\cite{lakonishok1992impact}에서 오나, 해당 문헌은 기관 \emph{내부}의 군집을,
식~\eqref{eq:herd}은 \emph{유형 간} 관계를 측정한다.

\textbf{(F3) 평균단가 대비 손실구간} 평균단가 대용치 $\bar P_{i,t}$와 당일 총거래
VWAP $P^{vwap}_{i,t}$로
\begin{equation}
x^{uw}_{i,t}=\max\!\left(
\frac{\bar P_{i,t}-P^{vwap}_{i,t}}{\bar P_{i,t}},\,0\right).
\label{eq:underwater}
\end{equation}
0에서 절단하므로 이 변수는 준거점 의존적 설계다. 평균단가는 총매수ㆍ총매도 수량과
금액으로 유효 보유량을 재귀 갱신해 구하되 전일 보유량에 망각계수 $\rho=0.98$을
곱한다(반감기 약 34거래일). 부호 기대는 처분효과
문헌~\cite{shefrin1985disposition,odean1998investors}에서 온다. 다만 이 대용치는
집계 거래로 만든 근사치이므로 $x^{uw}$의 계수를 처분효과의 직접 검정으로 읽을 수
없다.

\textbf{컨텍스트}는 KOSPI200 1일 수익률 $r_t$와 원/달러 환율 수준
$z_t=(\log e_t-\mu_{t-1,252})/\sigma_{t-1,252}$이다. 평균과 표준편차는 $t-1$까지의
252거래일만 쓴다. 삼성전자의 KOSPI200 비중 때문에 종목 모멘텀과 시장수익률이
겹치며 특징--컨텍스트 상관은 $-0.29$에서 $0.22$ 범위다. 따라서 개별 계수는 다른
항이 통제된 조건부 연관이다.

\subsection{근시안적 보상과 Lasso로의 환원}
\label{sec:model}

유형 $i$의 행동점수를
\begin{equation}
q_{i,t}=(\beta_i+B_iC_t)^{\top}x_{i,t}+\alpha_i^{\top}C_t
\label{eq:score}
\end{equation}
로 둔다. $\alpha$는 종목 특징이 평균일 때도 행동 수준을 옮기는 수준효과, $B$는
컨텍스트에 따라 특징 민감도가 달라지는 기울기 효과다. 상호작용을 넣으면서 하위차수
항 $C_t$를 빼는 것은 표준적 사양 오류이므로 $\alpha$의 포함은 자료로 정당화할
대상이 아니다. 연속 행동에 대한 단기 보상을
$R_i=q_{i,t}a-\tfrac{\kappa}{2}a^{2}$로 두면 최적 정책은
\begin{equation}
a^{*}_{i,t+1}=\mathrm{clip}(q_{i,t}/\kappa,-1,1)
\label{eq:policy}
\end{equation}
이다. 이차항이 없으면 해가 거의 항상 $\pm1$이 되어 행동 강도를 재현할 수 없다.
$\kappa$는 식별을 위해 1로 고정하며, 이는 보상과 계수의 단위를 고정하는 정규화이지
곡률이 경험적으로 1이라는 주장이 아니다. 학습목적은
\begin{equation}
\min_{\theta_i}\ \frac{1}{|\mathcal{D}_i|}
\sum_{t\in\mathcal{D}_i}\!\bigl(a^{*}_{i,t+1}-a_{i,t+1}\bigr)^{2}
+\lambda_i\lVert\theta_i\rVert_1
\label{eq:loss}
\end{equation}
이고 Adam~\cite{kingma2015adam}으로 유형별 독립 학습한다.

\begin{proposition}[Lasso 환원]
\label{prop:lasso}
$z_{i,t}=[\,x_{i,t};\ \mathrm{vec}(x_{i,t}C_t^{\top});\ C_t\,]\in\mathbb{R}^{11}$,
$\theta_i=[\,\beta_i;\ \mathrm{vec}(B_i);\ \alpha_i\,]$로 두면
$q_{i,t}=\theta_i^{\top}z_{i,t}$이다. 학습 표본 전체에서 $|q_{i,t}|\le 1$이면
식~\eqref{eq:policy}의 clip이 작동하지 않아 $a^{*}=q$이고, 식~\eqref{eq:loss}은
$a_{i,t+1}$을 $z_{i,t}$에 회귀하는 $L_1$ 회귀(Lasso)와 정확히 같아진다.
\end{proposition}

본 표본에서 $|q|$의 최대값은 외국인 0.842, 개인 0.587, 기관 0.135이고 45개 분할
전부에서 포화 비율이 정확히 0이므로, 명제~\ref{prop:lasso}의 전제가 실제로
성립한다. 즉 SPR-IRL에서 ``선호를 복원한다''와 ``수급을 회귀한다''는 수치적으로
구별되지 않는다.

이 일치는 결함이 아니라 근시안 가정의 귀결이다. 한 걸음만 내다보면 오늘의 행동이
내일 이후에 영향을 주지 않으므로, 무엇을 원하는가(보상)와 무엇을 하는가(정책)가
하나로 겹친다. 최대엔트로피 IRL~\cite{ziebart2008maximum}의
$\pi\propto\exp(R/\tau)$를 적용하면 $\pi$는 평균 $q$의 절단정규분포가 되어
점추정치가 위와 일치하므로, 이 환원은 특정 구현의 우연이 아니라 근시안ㆍ이차보상
설정 전반의 성질이다. 역강화학습 정식화의 가치는 확장 경로에 있다. 상태 동역학과
다기간 가치를 넣는 순간 보상은 정책과 갈라지고 회귀로는 도달할 수 없는 구조적
복원이 된다(\S\ref{sec:discussion}).

\subsection{검증 프로토콜}
\label{sec:protocol}

복원된 계수를 그대로 보고하지 않고 네 단계를 거친다. \textbf{(V1)} 10개 시간 폴드
중 2개를 시험으로 쓰는 조합형 정화 교차검증~\cite{lopez2018advances}으로
$\binom{10}{2}=45$개 분할을 만들고, 1일 purge와 5일 embargo를 적용하며 모든
표준화는 학습 인덱스에만 적합한다. 분할들이 관측을 공유하므로 부호 일관성은
$p$-값이 아니라 안정성 지표로만 해석한다. \textbf{(V2)} 그 한계를 보완하기 위해
표본을 달력월 블록으로 나누어 1{,}000회 복원추출하고 95\% 신뢰구간을 구한다. 월
블록은 일별 수급의 자기상관을 블록 안에 보존한다. \textbf{(V3)} CPCV는 시험 폴드가
학습 폴드보다 앞설 수 있으므로, 최초 400개 관측에서 시작해 5일 purge 후 다음
60거래일을 예측하고 학습 구간을 60거래일씩 넓히는 expanding walk-forward를 10개 창
수행한다. \textbf{(V4)} 결과가 $L_1$ 벌점의 산물이 아님을 보이기 위해 $L_2$ 및
무벌점으로 재추정해 부호를 비교한다.

지표는 MAE, RMSE, 방향 정확도($\mathrm{sign}(a^{*})=\mathrm{sign}(a)$의 비율),
예측--실제 상관, 포화 비율이다. 유형 간 행동 표준편차가 다르므로 오차의 절대값은
유형 간 비교에 쓰지 않는다. 자명한 참조로 전일 자기행동을 그대로 예측하는 지속성
규칙을 함께 보고한다.

% =====================================================================
\section{실험}
\label{sec:experiments}

\subsection{데이터}
\label{sec:data}

삼성전자(005930) 일별 자료에서 252일 환율 이력, 20일 모멘텀, 다음 날 행동을 모두
계산할 수 있는 최신 973행을 사용한다. 상태일은 2022-01-06--2025-12-29이다. 단일
종목은 표본의 한계가 아니라 종목 고정효과와 유동성 차이를 배제하고 유형 간 선호
이질성만 분리하기 위한 통제된 설정이다.

행동 평균은 세 유형 모두 0에 가깝지만 표준편차는 개인 0.334, 외국인 0.257, 기관
0.125로 다르고, 1차 자기상관은 각각 0.345, 0.401, 0.149다. 동일 거래일 유형 간
행동은 강한 음의 관계를 갖는다(외국인--개인 $-0.772$, 기관--개인 $-0.499$). 이는 한
유형의 매수가 다른 유형의 매도로 나타날 수밖에 없는 시장 청산 제약이며,
\S\ref{sec:beta}의 해석 범위를 제한한다.

학습 설정은 유형별로 외국인 75 epochs / batch 256 / lr 0.001 / $\lambda=0$, 기관
10 / 2048 / 0.0005 / 0.01, 개인 20 / 512 / 0.001 / 0.0003이며 seed 42, $\kappa=1$,
유형당 파라미터 11개($\beta$ 3, $B$ 6, $\alpha$ 2)다.

\subsection{표본외 행동 재현}
\label{sec:oos}

\begin{figure}[t]
  \centering
  \includegraphics[width=\columnwidth]{figures/fig_actual_vs_prediction.pdf}
  \caption{실제 행동과 CPCV 표본외 예측(20거래일 이동평균). 각 날짜의 9개 시험예측을
  먼저 평균했다. 괄호 안은 평활 계열 간 상관으로,
  표~\ref{tab:performance}의 일별 상관과 직접 비교하지 않는다.}
  \Description{외국인, 기관, 개인의 실제 순매수와 모형 예측을 시간축에 겹쳐 그린
  세 개의 선그래프. 기관 패널에서 예측선이 0 근방에 평탄하게 놓여 있다.}
  \label{fig:prediction}
\end{figure}

\begin{table}[t]
  \caption{CPCV 표본외 성능(45개 분할 평균 $\pm$ 표준편차)과 지속성 기준선}
  \label{tab:performance}
  \small
  \centering
  \begin{tabular}{@{}lccc@{}}
    \toprule
    지표 & 외국인 & 기관 & 개인 \\
    \midrule
    \multicolumn{4}{@{}l}{\emph{SPR-IRL}}\\
    방향 정확도 & $0.636\pm0.039$ & $0.542\pm0.038$ & $0.608\pm0.044$ \\
    상관계수    & $0.278\pm0.110$ & $0.068\pm0.064$ & $0.245\pm0.118$ \\
    MAE         & $0.195\pm0.019$ & $0.094\pm0.008$ & $0.264\pm0.029$ \\
    MAE / 행동 SD & $0.759$ & $0.756$ & $0.791$ \\
    포화 비율   & $0.000$ & $0.000$ & $0.000$ \\
    \addlinespace
    \multicolumn{4}{@{}l}{\emph{지속성 기준선} ($a_{t+1}\!\leftarrow\!a_t$)}\\
    방향 정확도 & $\mathbf{0.642}$ & $\mathbf{0.547}$ & $\mathbf{0.616}$ \\
    상관계수    & $\mathbf{0.359}$ & $\mathbf{0.136}$ & $\mathbf{0.311}$ \\
    \bottomrule
  \end{tabular}
\end{table}

그림~\ref{fig:prediction}에서 외국인과 개인의 예측 계열은 실제 계열의 국면 전환을
따라가지만(평활 상관 0.69, 0.76) 진폭은 $L_1$ 축소 때문에 작다. 기관의 예측 계열은
사실상 0선에 붙어 있다(0.02).

표~\ref{tab:performance}에서 두 가지가 드러난다. 첫째, 분할 간 산포가 커서 단일
분할의 수치로 유형을 비교할 수 없다. 둘째, \textbf{자명한 지속성 기준선이 방향
정확도와 상관계수에서 SPR-IRL을 앞선다}. $a_t$는 $t$일 종가 이후 공개되므로 이는
정보상 이용 가능한 규칙이다. 원인은 명확하다. 행동의 자기상관이
$0.401/0.149/0.345$인데 보상 특징 집합에 자기 지연행동이 없다. 특징을
(F1)--(F3)로 제한한 것은 세 유형이 독립적으로 확립된 규칙성에 대응하는 특징만
공유하도록 한 설계이지 예측 성능을 위한 설계가 아니다. 표~\ref{tab:performance}은
모형이 잡음만 학습한 것이 아님을 보이는 용도이며 예측 우위의 근거가 아니다.

모든 유형ㆍ모든 분할의 포화 비율이 0이므로 명제~\ref{prop:lasso}의 전제가 성립하고,
SPR-IRL의 추정은 문자 그대로 Lasso다.

\subsection{복원된 선호: 모멘텀 축의 거울상}
\label{sec:beta}

\begin{figure*}[t]
  \centering
  \includegraphics[width=\textwidth]{figures/fig_beta_alpha_forest.pdf}
  \caption{복원된 기본 보상가중치 $\beta$(A)와 컨텍스트 주효과 $\alpha$(B)의 안정성.
  회색 점은 45개 개별 CPCV 분할, 굵은 오차선은 분할 간 표준편차, 옅은 가로선은 달력월
  블록 부트스트랩 95\% 신뢰구간이다. 속이 빈 표식은 신뢰구간이 0을 포함하는 항이다.}
  \Description{두 패널의 forest plot. 패널 A는 세 유형의 모멘텀, 군집, 손실구간
  계수를, 패널 B는 KOSPI200과 환율 주효과를 보인다. 외국인 모멘텀은 양, 개인
  모멘텀은 음이며 기관의 모든 표식은 속이 비어 있다.}
  \label{fig:forest}
\end{figure*}

\textbf{모멘텀 축의 거울상.} 외국인의 모멘텀 가중치는 $+0.0496$으로 45개 분할
전부에서 양수, 개인은 $-0.0304$로 전부에서 음수이며 두 부트스트랩 신뢰구간은 서로
겹치지 않고 각각 0을 제외한다(그림~\ref{fig:forest}A). 동일한 특징 집합을 유형별로
독립 추정했는데도 두 유형이 정반대로 갈리며, 이는 선행
연구~\cite{grinblatt2000investment,choe1999foreign}의 부호 기대와 일치한다.

\textbf{컨텍스트에서의 재현.} 같은 거울상이 주효과에서도 나타난다
(그림~\ref{fig:forest}B). 외국인은 시장 상승일에 순매수 방향으로($+0.0361$), 원화
약세 국면에 순매도 방향으로($-0.0261$) 이동하고 개인은 두 컨텍스트 모두에서 정확히
반대다($-0.0246$, $+0.0298$). 네 계수 모두 45개 분할 전부에서 부호가 유지되고
신뢰구간도 0을 제외한다.

\textbf{군집.} 군집 계수는 세 유형 모두 45개 분할 전부에서 음수이고 외국인
($-0.0402$)과 개인($-0.0322$)은 신뢰구간도 0을 제외한다. 그러나 이 결과는 행동적
해석과 시장 청산 제약을 분리하지 못한다. \S\ref{sec:data}의 $-0.77$이라는 동시점
상관 아래에서는 음의 군집 계수 중 상당 부분이 회계적 제약에서 발생할 수 있으므로,
본 연구는 이를 ``관측된 집계 수급이 유형 간 음의 교차 연관을 보인다''는 서술적
결과로 제한한다. 거울상 주효과에도 같은 유보가 적용된다.

\textbf{손실구간.} 개인의 손실구간 가중치는 $+0.0293$으로 45개 분할 전부에서
양수이며 처분효과 문헌~\cite{shefrin1985disposition,odean1998investors}의 기준가격
의존성과 정합한다. 그러나 그림~\ref{fig:forest}A에서 보듯 세 유형 모두 신뢰구간이
0을 포함한다. 개인의 구간은 $[-0.0122,+0.0665]$로, 45개 분할에서 부호가 100\%
일관됨에도 월 블록 재표집에서는 5.5\%가 부호를 뒤집는다. 이는 CPCV 분할이 관측을
공유하기 때문에 부호 일관성이 불확실성을 과소평가한다는 \S\ref{sec:protocol}의
경고가 실제로 작동한 사례다. 따라서 처분효과와 정합하는 부호를 관측했다고 보고하되
통계적으로 확립된 결과로 주장하지 않는다.

상호작용 $B$ 12개 중 45개 분할 전부에서 부호가 유지된 것은 외국인
모멘텀$\times$KOSPI200($-0.0139$), 개인 모멘텀$\times$환율($+0.0144$), 개인
군집$\times$KOSPI200($-0.0080$), 개인 손실구간$\times$환율($-0.0186$) 넷이고, 이 중
부트스트랩을 통과하는 것은 첫째와 넷째다. 12개를 동시에 검토했으므로 상호작용은
탐색적 결과로만 제시한다.

\subsection{안정성}
\label{sec:stability}

\begin{figure*}[t]
  \centering
  \includegraphics[width=\textwidth]{figures/fig_walk_forward.pdf}
  \caption{Expanding walk-forward 10개 창의 계수 경로와 방향 정확도. 성능은 창마다
  흔들리지만 계수 부호는 유지되며, 외국인과 개인의 모멘텀 경로는 한 번도 교차하지
  않는다.}
  \Description{확장 창 번호를 가로축으로 하여 세 유형의 모멘텀 계수, 두 컨텍스트
  주효과, 방향 정확도를 그린 네 개의 선그래프.}
  \label{fig:walkforward}
\end{figure*}

\textbf{(V2)} 부트스트랩 결과는 셋으로 갈린다. 모멘텀 거울상과 컨텍스트 거울상을
이루는 8개 항이 모두 0을 제외하고, 손실구간은 통과하지 못하며, 기관은 어떤 항도
0을 제외하지 못한다(그림~\ref{fig:forest}의 속 빈 표식).

\textbf{(V3)} 그림~\ref{fig:walkforward}에서 성능은 크게 흔들린다. 외국인 방향
정확도는 0.43--0.85를 오가며 평균 $0.610\pm0.116$, 상관은 $0.131\pm0.162$로 CPCV의
0.278에서 절반 이하다. CPCV는 시험 폴드를 학습 폴드보다 앞세울 수 있는 반면
walk-forward는 그렇지 않으므로, 본 연구는 이 저하를 시간 구조에서 오는 일반화
한계로 그대로 보고한다. 반면 계수 경로는 흔들리지 않는다. 외국인 모멘텀은 10개 창
전부 양수, 개인은 전부 음수이며 두 경로가 한 번도 교차하지 않고, 두 컨텍스트
주효과에서도 같은 거울상이 유지된다. 기관의 경로는 0선에 붙어 있다.

\textbf{(V4)} 외국인과 개인의 계수는 $L_1$, $L_2$, 무벌점에서 소수점 넷째 자리까지
일치하므로 거울상은 벌점의 산물이 아니다. 반면 기관의 계수만 벌점에 따라
이동한다(군집 $-0.0076\rightarrow-0.0124$, 손실구간
$-0.0000\rightarrow-0.0018$).

\subsection{기관: 수렴된 null}
\label{sec:null}

기관 수급에 대해 SPR-IRL은 일별 빈도에서 안정적으로 식별되는 선호를 찾지 못했고,
본 연구는 이를 은폐하지 않고 결과로 보고한다. 근거는 다섯이다. (i) 예측 계열이 0선에
퇴화한다(그림~\ref{fig:prediction}). (ii) 군집을 제외한 계수가 다른 유형의 $1/40$
수준이다($-0.0006$, $-0.0008$). (iii) 부트스트랩에서 0을 제외하는 항이
없다(그림~\ref{fig:forest}). (iv) 계수가 벌점 선택에 좌우된다. (v) walk-forward
10개 창에서 모멘텀 부호 일관성이 40\%, 손실구간이 0\%다.

\textbf{이 null은 최적화 실패가 아니다.} 명제~\ref{prop:lasso}에 따라
식~\eqref{eq:loss}은 볼록 문제이므로 전역 최적해가 유일하다. 도달한 해와 동일
목적함수의 폐쇄형 최적해를 비교하면 기관의 목적함수 상대 격차는 평균 0.17\%, 최대
0.48\%에 불과하다(외국인 0.21\%, 개인 0.92\%). 세 유형 모두 도달 가능한 최적해에
근접했고, 기관에서만 그 최적해의 계수가 0 근방이다. 데이터가 기관에 대해 결정하는
선호가 없는 것이지 최적화가 찾지 못한 것이 아니다.

경제적으로 이 null은 자연스럽다. ``기관''은 연기금, 투신, 보험, 사모, 금융투자를
합산한 범주이므로 서로 다른 목적함수를 가진 주체들의 수급이 상쇄된다.

\subsection{구성 타당성 점검}
\label{sec:construct}

\begin{table}[t]
  \caption{사전 부호 기대와 복원 결과}
  \label{tab:construct}
  \small
  \begin{tabular}{@{}>{\raggedright\arraybackslash}p{0.30\columnwidth}
                  >{\raggedright\arraybackslash}p{0.13\columnwidth}
                  >{\raggedright\arraybackslash}p{0.26\columnwidth}
                  >{\raggedright\arraybackslash}p{0.16\columnwidth}@{}}
    \toprule
    규칙성 & 사전 기대 & 복원 결과 & 판정 \\
    \midrule
    외국인 추세거래~\cite{grinblatt2000investment,choe1999foreign}
      & $\beta^{mom}\!>\!0$ & $+0.0496$ (100\%, CI 0 제외) & 부합 \\
    개인 역추세~\cite{grinblatt2000investment,kaniel2008individual}
      & $\beta^{mom}\!<\!0$ & $-0.0304$ (100\%, CI 0 제외) & 부합 \\
    처분효과~\cite{shefrin1985disposition,odean1998investors}
      & $\beta^{uw}\!>\!0$ & $+0.0293$ (100\%, CI 0 포함) & 부분 부합 \\
    기관 군집행동~\cite{lakonishok1992impact}
      & $\beta^{herd}\!>\!0$ & $-0.0044$ (100\%, CI 0 포함) & 불일치 \\
    \bottomrule
  \end{tabular}
\end{table}

표~\ref{tab:construct}이 대조 결과다. 동일한 특징과 모형 형태를 부여하고 파라미터만
분리 학습했음에도 복원된 부호가 독립적으로 확립된 규칙성과 부합하고 두 유형이
모멘텀에서 정반대로 갈린다는 점이 구성 타당성의 증거다. 불일치 항목은 측정 대상이
다르기 때문이다. Lakonishok 외~\cite{lakonishok1992impact}는 기관 \emph{내부}의
군집을, 식~\eqref{eq:herd}은 기관이 \emph{다른 두 유형}과 같은 방향으로 움직이는지를
측정하므로 본 결과는 해당 문헌의 반증이 아니다.

% =====================================================================
%  5장 논의와 한계 (한국어)
%
%  '본 연구가 주장하지 않는 것'은 4장(결과)이 아니라 이 절에 둔다.
%  분량 설계: 약 0.6쪽.
% =====================================================================

\section{논의와 한계}
\label{sec:discussion}

SPR-IRL이 복원한 것은 세 유형의 선호를 하나의 공통 척도 위에 나란히 놓은 그림이다.
모멘텀 축에서 외국인과 개인이 정반대로 갈리고 같은 대비가 컨텍스트 반응에서
재현되며, 기관은 군집을 제외하면 식별되는 선호를 보이지 않는다. 이 세 결과는 세
유형에 동일한 특징과 동일한 모형 형태를 부여한 뒤 파라미터만 분리 학습해 얻은
것이므로 설계가 심어둔 대비가 아니다.

\subsection{본 연구가 주장하지 않는 것}
\label{sec:notclaim}

프로토콜의 마지막 단계로, 결과가 지지하지 않는 주장을 명시한다.

\textbf{(1) 수익률 예측가능성.} 본 연구의 지표는 수급의 방향과 강도에 대한 재현
오차이며 수익률, 샤프지수, 거래비용을 포함하지 않는다. 표~\ref{tab:performance}에서
어떤 매매 규칙도 도출되지 않는다.

\textbf{(2) 수급 예측 성능의 우위.} 지속성 기준선이 방향 정확도와 상관계수에서
SPR-IRL을 앞선다(\S\ref{sec:oos}). 선호 복원과 예측은 다른 특징 집합을 요구한다.
예비 점검에서 지연 자기행동을 네 번째 특징으로 추가하면 모형이 기준선을
넘어서지만(방향 정확도 $0.646/0.569/0.629$) 개인의 $\alpha^{KOSPI}$ 부호가 반전되고
손실구간 계수가 소멸한다. 즉 \textbf{본 논문이 보고한 거울상 중 일부는 수급
자기상관과 분리되지 않는다.}

\textbf{(3) 인과효과.} 복원된 가중치는 표준화된 관측 특징 사이의 조건부 연관이다.
특징--컨텍스트 상관이 $-0.29$에서 $0.22$ 범위이므로 개별 계수의 크기는 다른 항의
통제 조건에 의존한다.

\textbf{(4) 유형 간 음의 연관의 행동적 해석.} 군집 계수와 거울상 컨텍스트 반응은
시장 청산 제약과 분리되지 않는다(\S\ref{sec:beta}).

\textbf{(5) 처분효과의 검정.} 개인의 손실구간 부호는 문헌과 정합하나 월 블록
부트스트랩을 통과하지 못했고, 평균단가 대용치는 계좌별 실현손익이 아니라 집계
거래로 만든 근사치다.

\textbf{(6) 종목 간 일반화.} 단일 종목 설정은 유형 간 이질성 분리를 위한 통제이며
다른 종목과 시장으로의 확장은 검증되지 않았다.

\subsection{향후 과제}

가장 시급한 것은 지연 자기행동을 특징으로 편입해 자기상관을 통제한 뒤 거울상이
살아남는지 확인하는 일이다. 다음으로 ``기관''을 연기금ㆍ투신ㆍ보험 등 세부 주체로
분해하면 상쇄로 가려진 선호가 드러날 수 있다. 마지막으로 본 논문의 근시안 가정을
푸는 것이 남는다. 명제~\ref{prop:lasso}가 보이듯 한 걸음만 내다보는 설정에서 보상은
정책과 겹치지만, 상태 동역학과 다기간 가치를 도입하면 두 대상이 갈라지고 회귀로는
도달할 수 없는 구조적 복원이 된다. SPR-IRL의 대칭 추정과 검증 프로토콜은 그
확장에서도 그대로 쓰일 수 있다.


\bibliographystyle{ACM-Reference-Format}
\bibliography{paper-ko,paper-refs-additional}

\end{document}
